\documentclass[aspectratio=169]{beamer}
%--- Configuración del Tema y Colores ---
\usetheme{Madrid}
\usecolortheme{whale}
\setbeamertemplate{navigation symbols}{}
\setbeamertemplate{caption}[numbered]

%--- Paquetes ---
\usepackage[utf8]{inputenc}
\usepackage[spanish]{babel}
\usepackage{graphicx}
\usepackage{booktabs}
\usepackage{tikz}
\usepackage{listings}
\usepackage{xcolor}

%--- Configuración de Listings para Python ---
\definecolor{codegreen}{rgb}{0,0.6,0}
\definecolor{codegray}{rgb}{0.5,0.5,0.5}
\definecolor{codepurple}{rgb}{0.58,0,0.82}
\definecolor{backcolour}{rgb}{0.95,0.95,0.92}

\lstdefinestyle{pythonstyle}{
    backgroundcolor=\color{backcolour},
    commentstyle=\color{codegreen},
    keywordstyle=\color{blue}\bfseries,
    numberstyle=\tiny\color{codegray},
    stringstyle=\color{codepurple},
    basicstyle=\ttfamily\scriptsize,
    breakatwhitespace=false,
    breaklines=true,
    captionpos=b,
    keepspaces=true,
    numbers=left,
    numbersep=5pt,
    showspaces=false,
    showstringspaces=false,
    showtabs=false,
    tabsize=2,
    language=Python,
    morekeywords={lambda, map, filter, reduce, True, False}
}
\lstset{style=pythonstyle}

%--- Metadatos del Documento ---
\title[PM2]{Programación Matemática 2}
\subtitle{Paradigma Imperativo y Funcional}
\author[Luis Alfredo Alvarado]{Mgtr. Luis Alfredo Alvarado Rodríguez}
\institute[ECFM - USAC]{
    Escuela de Ciencias Físicas y Matemáticas \\
    Universidad de San Carlos de Guatemala
}
\date{Primer Semestre, 2026}

\begin{document}

%--- Diapositiva de Título ---
\begin{frame}
    \titlepage
\end{frame}

%--- Diapositiva: Tabla de Contenidos ---
\begin{frame}{Agenda de Hoy}
    \tableofcontents
\end{frame}

%===========================================================
% SECCIÓN 1: INTRODUCCIÓN
%===========================================================
\section{¿Qué es un Paradigma de Programación?}

\begin{frame}{Paradigmas de Programación}
    \begin{block}{Definición}
        Un \textbf{paradigma de programación} es una forma fundamental de pensar y estructurar el código. Define \textit{cómo} se describe la solución a un problema.
    \end{block}
    \vspace{0.5cm}
    
    \textbf{Principales paradigmas:}
    \begin{itemize}
        \item \textbf{Imperativo:} Se describe \textit{cómo} hacer algo (secuencia de instrucciones).
        \item \textbf{Funcional:} Se describe \textit{qué} resultado se desea (transformaciones).
        \item \textbf{Orientado a Objetos:} Se modela el mundo con objetos y sus interacciones.
        \item Declarativo, Lógico, Reactivo...
    \end{itemize}
    \vspace{0.3cm}
    \centering
    \textit{Python es un lenguaje \textbf{multiparadigma}: soporta todos estos estilos.}
\end{frame}

%===========================================================
% SECCIÓN 2: PARADIGMA IMPERATIVO
%===========================================================
\section{Paradigma Imperativo}

\begin{frame}{Paradigma Imperativo: Concepto}
    \begin{alertblock}{Filosofía}
        \textbf{``Se proporcionan instrucciones paso a paso de CÓMO resolver el problema.''}
    \end{alertblock}
    \vspace{0.3cm}
    
    \textbf{Características:}
    \begin{itemize}
        \item El programa es una \textbf{secuencia de comandos} que modifican el estado.
        \item Se hace uso extensivo de \textbf{variables mutables}.
        \item Estructuras de control: \texttt{if}, \texttt{for}, \texttt{while}.
        \item El orden de las instrucciones \textbf{importa}.
    \end{itemize}
    \vspace{0.3cm}
    
    \textbf{Analogía:} Una receta de cocina con pasos numerados.
\end{frame}

\begin{frame}[fragile]{Ejemplo Imperativo: Suma de una Lista}
    \textbf{Problema:} Calcular la suma de los números del 1 al 10.
    \vspace{0.3cm}
    
    \begin{lstlisting}[title=Enfoque Imperativo]
# Inicializar estado (variable mutable)
suma = 0

# Secuencia de instrucciones que modifican el estado
for i in range(1, 11):
    suma = suma + i  # Se muta 'suma' en cada iteracion

print(f"La suma es: {suma}")
# Salida: La suma es: 55
    \end{lstlisting}
    
    \vspace{0.3cm}
    \textbf{Nótese:} La variable \texttt{suma} cambia 10 veces durante la ejecución.
\end{frame}

\begin{frame}[fragile]{Ejemplo Imperativo: Filtrar Números Pares}
    \textbf{Problema:} Obtener los números pares de una lista.
    \vspace{0.3cm}
    
    \begin{lstlisting}[title=Enfoque Imperativo]
numeros = [1, 2, 3, 4, 5, 6, 7, 8, 9, 10]
pares = []  # Lista mutable vacia

for num in numeros:
    if num % 2 == 0:  # Condicion
        pares.append(num)  # Se modifica la lista

print(pares)
# Salida: [2, 4, 6, 8, 10]
    \end{lstlisting}
    
    \vspace{0.3cm}
    \textbf{Patrón típico:} Inicializar $\rightarrow$ Iterar $\rightarrow$ Mutar $\rightarrow$ Retornar
\end{frame}

\begin{frame}[fragile]{Ejemplo Imperativo: Factorial}
    \textbf{Problema:} Calcular $n! = n \times (n-1) \times \ldots \times 1$
    \vspace{0.3cm}
    
    \begin{lstlisting}[title=Factorial Imperativo]
def factorial_imperativo(n):
    resultado = 1  # Estado inicial
    
    for i in range(1, n + 1):
        resultado = resultado * i  # Mutacion
    
    return resultado

print(factorial_imperativo(5))
# Salida: 120
    \end{lstlisting}
    
    \vspace{0.3cm}
    \textbf{Flujo mental:} $1 \rightarrow 1 \rightarrow 2 \rightarrow 6 \rightarrow 24 \rightarrow 120$
\end{frame}

%===========================================================
% SECCIÓN 3: PARADIGMA FUNCIONAL
%===========================================================
\section{Paradigma Funcional}

\begin{frame}{Paradigma Funcional: Concepto}
    \begin{exampleblock}{Filosofía}
        \textbf{``Se especifica QUÉ resultado se desea, no cómo calcularlo paso a paso.''}
    \end{exampleblock}
    \vspace{0.3cm}
    
    \textbf{Características:}
    \begin{itemize}
        \item Las funciones son \textbf{ciudadanos de primera clase} (se pasan como argumentos).
        \item \textbf{Inmutabilidad:} Se evita cambiar el estado de las variables.
        \item \textbf{Funciones puras:} Mismo input $\Rightarrow$ mismo output (sin efectos secundarios).
        \item Uso de \texttt{map}, \texttt{filter}, \texttt{reduce}, \texttt{lambda}.
    \end{itemize}
    \vspace{0.3cm}
    
    \textbf{Analogía:} Una función matemática $f(x) = x^2$ siempre da el mismo resultado.
\end{frame}

\begin{frame}[fragile]{Funciones Lambda (Anónimas)}
    \begin{block}{Sintaxis}
        \texttt{lambda argumentos: expresión}
    \end{block}
    \vspace{0.3cm}
    
    \begin{lstlisting}[title=Ejemplos de Lambda]
# Funcion tradicional
def cuadrado(x):
    return x ** 2

# Equivalente con lambda
cuadrado_lambda = lambda x: x ** 2

print(cuadrado(5))         # 25
print(cuadrado_lambda(5))  # 25

# Lambda con multiples argumentos
suma = lambda a, b: a + b
print(suma(3, 4))  # 7
    \end{lstlisting}
\end{frame}

\begin{frame}[fragile]{Ejemplo Funcional: Suma con \texttt{reduce}}
    \textbf{Problema:} Calcular la suma de los números del 1 al 10.
    \vspace{0.3cm}
    
    \begin{lstlisting}[title=Enfoque Funcional]
from functools import reduce

numeros = range(1, 11)

# reduce(funcion, iterable) -> aplica funcion acumulativamente
suma = reduce(lambda acc, x: acc + x, numeros)

print(f"La suma es: {suma}")
# Salida: La suma es: 55

# Alternativa aun mas funcional (funcion built-in)
suma_directa = sum(range(1, 11))
    \end{lstlisting}
    
    \vspace{0.3cm}
    \textbf{Nótese:} No existen variables que muten. Es una \textit{transformación} de datos.
\end{frame}

\begin{frame}[fragile]{Ejemplo Funcional: Filtrar con \texttt{filter}}
    \textbf{Problema:} Obtener los números pares de una lista.
    \vspace{0.3cm}
    
    \begin{lstlisting}[title=Enfoque Funcional]
numeros = [1, 2, 3, 4, 5, 6, 7, 8, 9, 10]

# filter(funcion_booleana, iterable)
es_par = lambda x: x % 2 == 0
pares = list(filter(es_par, numeros))

print(pares)
# Salida: [2, 4, 6, 8, 10]

# Version con list comprehension (estilo Pythonic)
pares_lc = [x for x in numeros if x % 2 == 0]
    \end{lstlisting}
    
    \vspace{0.3cm}
    \textbf{Nota:} \texttt{filter} retorna un iterador, por eso se utiliza \texttt{list()}.
\end{frame}

\begin{frame}[fragile]{Ejemplo Funcional: Transformar con \texttt{map}}
    \textbf{Problema:} Elevar al cuadrado todos los números de una lista.
    \vspace{0.3cm}
    
    \begin{lstlisting}[title=Enfoque Funcional con map]
numeros = [1, 2, 3, 4, 5]

# map(funcion, iterable) -> aplica funcion a cada elemento
cuadrados = list(map(lambda x: x ** 2, numeros))

print(cuadrados)
# Salida: [1, 4, 9, 16, 25]

# Equivalente con list comprehension
cuadrados_lc = [x ** 2 for x in numeros]
    \end{lstlisting}
    
    \vspace{0.3cm}
    \textbf{Concepto clave:} \texttt{map} transforma \textit{cada} elemento sin modificar la lista original.
\end{frame}

\begin{frame}[fragile]{Ejemplo Funcional: Factorial Recursivo}
    \textbf{Problema:} Calcular $n!$ sin usar ciclos ni variables mutables.
    \vspace{0.3cm}
    
    \begin{lstlisting}[title=Factorial Funcional (Recursivo)]
def factorial_funcional(n):
    # Caso base
    if n <= 1:
        return 1
    # Caso recursivo: n! = n * (n-1)!
    return n * factorial_funcional(n - 1)

print(factorial_funcional(5))
# Salida: 120
    \end{lstlisting}
    
    \vspace{0.3cm}
    \textbf{Flujo:} $5! = 5 \times 4! = 5 \times 4 \times 3! = \ldots = 5 \times 4 \times 3 \times 2 \times 1$
    
    \vspace{0.2cm}
    \small\textit{La recursión es el ``loop'' del paradigma funcional.}
\end{frame}

\begin{frame}[fragile]{Composición de Funciones}
    \textbf{Problema:} Obtener la suma de los cuadrados de los números pares.
    \vspace{0.3cm}
    
    \begin{lstlisting}[title=Pipeline Funcional]
from functools import reduce

numeros = [1, 2, 3, 4, 5, 6, 7, 8, 9, 10]

# Pipeline: filtrar -> transformar -> reducir
resultado = reduce(
    lambda acc, x: acc + x,          # 3. Sumar
    map(
        lambda x: x ** 2,            # 2. Elevar al cuadrado
        filter(lambda x: x % 2 == 0, numeros)  # 1. Filtrar pares
    )
)

print(resultado)  # 2^2 + 4^2 + 6^2 + 8^2 + 10^2 = 220
    \end{lstlisting}
    
    \vspace{0.2cm}
    \textbf{Lectura:} ``Suma de los cuadrados de los pares de números''
\end{frame}

%===========================================================
% SECCIÓN 4: COMPARACIÓN
%===========================================================
\section{Comparación: Imperativo vs Funcional}

\begin{frame}{Comparación Directa}
    \begin{table}[]
        \centering
        \renewcommand{\arraystretch}{1.4}
        \begin{tabular}{p{4cm} p{5cm} p{5cm}}
            \toprule
            \textbf{Aspecto} & \textbf{Imperativo} & \textbf{Funcional} \\
            \midrule
            Enfoque & \textit{Cómo} hacerlo & \textit{Qué} resultado \\
            Estado & Variables mutables & Inmutabilidad \\
            Control de flujo & \texttt{for}, \texttt{while}, \texttt{if} & Recursión, \texttt{map}, \texttt{filter} \\
            Efectos secundarios & Comunes & Evitados \\
            Depuración & Seguir paso a paso & Probar funciones aisladas \\
            Legibilidad & Más verboso & Más declarativo \\
            \bottomrule
        \end{tabular}
    \end{table}
\end{frame}

\begin{frame}[fragile]{Mismo Problema, Dos Estilos}
    \textbf{Problema:} Convertir lista de nombres a mayúsculas y filtrar los que empiezan con 'A'.
    
    \begin{columns}
        \column{0.48\textwidth}
        \begin{lstlisting}[title=Imperativo]
nombres = ["ana", "bob", 
           "alberto", "carlos"]
resultado = []

for nombre in nombres:
    mayus = nombre.upper()
    if mayus.startswith('A'):
        resultado.append(mayus)

print(resultado)
# ['ANA', 'ALBERTO']
        \end{lstlisting}
        
        \column{0.48\textwidth}
        \begin{lstlisting}[title=Funcional]
nombres = ["ana", "bob", 
           "alberto", "carlos"]

resultado = list(
    filter(
        lambda x: x.startswith('A'),
        map(str.upper, nombres)
    )
)

print(resultado)
# ['ANA', 'ALBERTO']
        \end{lstlisting}
    \end{columns}
\end{frame}

\begin{frame}{¿Cuándo Utilizar Cada Uno?}
    \begin{columns}[t]
        \column{0.48\textwidth}
        \begin{alertblock}{Se recomienda Imperativo cuando...}
            \begin{itemize}
                \item Se necesita control preciso del flujo.
                \item El rendimiento es crítico (menos overhead).
                \item El algoritmo es inherentemente secuencial.
                \item Se trabaja con I/O o efectos secundarios.
            \end{itemize}
        \end{alertblock}
        
        \column{0.48\textwidth}
        \begin{exampleblock}{Se recomienda Funcional cuando...}
            \begin{itemize}
                \item Se procesan colecciones de datos.
                \item Se busca código más conciso y declarativo.
                \item Se necesita paralelizar fácilmente.
                \item Se requieren funciones testeables y predecibles.
            \end{itemize}
        \end{exampleblock}
    \end{columns}
    
    \vspace{0.5cm}
    \centering
    \textbf{En Python:} Lo mejor de ambos mundos. Se recomienda mezclar según convenga.
\end{frame}

%===========================================================
% SECCIÓN 5: EJERCICIOS
%===========================================================
\section{Ejercicios Propuestos}

\begin{frame}{Ejercicios para Practicar}
    \textbf{Resolver cada problema en AMBOS estilos:}
    \vspace{0.5cm}
    
    \begin{enumerate}
        \setlength\itemsep{0.8em}
        \item Dada una lista de números, retornar solo los mayores a 10.
        \item Calcular el producto de todos los elementos de una lista.
        \item Dada una lista de strings, retornar la concatenación de todos.
        \item Implementar la función \texttt{fibonacci(n)} que retorne el n-ésimo número de Fibonacci.
        \item \textbf{(Bonus)} Dada una lista de diccionarios con estudiantes \texttt{\{"nombre": ..., "nota": ...\}}, obtener el promedio de notas de quienes aprobaron (nota $\geq$ 61).
    \end{enumerate}
\end{frame}

%--- Diapositiva Final ---
\begin{frame}
    \centering
    \Huge \textbf{¿Preguntas?}
    
    \vspace{0.8cm}
    \large
    \textbf{Próxima clase:} Paradigma Orientado a Objetos
    
    \vspace{0.5cm}
    \normalsize
    Mgtr. Luis Alfredo Alvarado Rodríguez\\
    \small Escuela de Ciencias Físicas y Matemáticas
    
    \vspace{0.5cm}
    \footnotesize
    \textit{``Primero se resuelve el problema. Después, se escribe el código.''} — John Johnson
\end{frame}

\end{document}